\section{Marco teórico y trabajos relacionados}
Se consultaron artículos académicos que sirven como referencia para justificar las 
decisiones tecnológicas y metodológicas. El estudio sobre Laravel [1] recalca la 
relevancia de este framework en términos de seguridad y rapidez, mientras que el trabajo 
comparativo entre Laravel y Symfony [7] permitió confirmar que, aunque Symfony 
ofrece mayor robustez empresarial, Laravel resulta más accesible y flexible para un 
equipo en formación. 

Del mismo modo, React se presenta como una biblioteca flexible y fácil de adoptar, lo 
que coincide con investigaciones que la destacan como la opción preferida frente a 
Angular y Vue en numerosos proyectos [2], [4], [5]. React Native, por su parte, 
constituye una herramienta estratégica para futuras ampliaciones hacia entornos 
multiplataforma [3]. 

Respecto a metodologías ágiles, el marco de trabajo Scrum se utilizó como base para 
organizar el flujo de trabajo, asegurando entregas incrementales y revisiones continuas 
[6]. Como lo advierte la literatura sobre metodologías ágiles [9], su adopción debe 
complementarse con buenas prácticas de ingeniería para asegurar escalabilidad. 